\chapter{Work log: what has been changed and pushed}
\label{ch:worklog}

Repository: \file{github.com/chalach49266/nv-compact-magnetometer}, branch
\co{main}. All commits below are pushed. The two sessions are separated because the
second one partly retracts the first.

\begin{retractbox}[title=\textbf{Read Round 7 first --- the acquisition path was reverted}]
On 2026-08-18 the acquisition \emph{modules} were reset to revision \co{8fd303f},
the one the second machine takes good data with. Everything Rounds 3, 4 and 6
changed on that path is \textbf{no longer in the code}. The notebooks were kept.
Rounds 1, 2 and 5 stand: measurement, analysis and figure saving, none of which
touch acquisition.
\end{retractbox}

\section{Round 1 --- 2026-08-12, against the 2026-08-06 session}

\begin{center}
\small
\begin{tabular}{@{}llp{8.6cm}@{}}
\toprule
\hd{Commit} & \hd{Scope} & \hd{What it did} \\
\midrule
\co{5ea75ab} & \co{data:} & 2026-08-06 session snapshot (drone sweeps, burst runs,
  17-point integration-time scan). \\
\co{50fe514} & \co{docs:} & The timing \& noise analysis, plus
  \co{analyze_twopoint_session.py} and \co{profile_twopoint_acquire.py}. \\
\co{942d9e8} & \co{perf:} & Corrected the \co{config_all} keyword; $\tau$ 213
  $\rightarrow$ \SI{120}{\micro\second}; added ABBA ordering; append-only CSV writer. \\
\co{6765284} & \co{fix:} & \co{MultipointLockinODMR.stream()} --- the continuous
  streaming path. \\
\co{9eed2e1} & \co{feat:} & \co{clean_burst_lockin.py}, to salvage burst runs already
  recorded. \\
\co{16fd9ba} & \co{docs:} & Changelog, rig checklist, cross-references. \\
\co{f700223} & \co{fix:} & Sensitivity omitted the $\sqrt{\text{reps}}$ factor,
  understating $\eta$ by $\sim$10$\times$. \\
\bottomrule
\end{tabular}
\end{center}

\paragraph{Scorecard on round 1.} Of the four changes intended to improve things:

\begin{center}
\small
\begin{tabular}{@{}p{6.0cm}lp{6.4cm}@{}}
\toprule
\hd{Change} & \hd{Outcome} & \hd{Verdict} \\
\midrule
$\tau$ 213 $\rightarrow$ \SI{120}{\micro\second} & \Partial &
  Took effect and is right for sensitivity, but bought \textbf{no rate} --- the
  premise that $\tau$ set the rate was wrong. \\
\co{config_all} fix + \co{reset_tproc} & \Open &
  Real bug, correctly fixed, but \textbf{did not stop the replay}. \\
\co{stream()} & \Done &
  Works: \textbf{1041.7\,Hz}, gapless. Its noise is a new open defect. \\
Append-only CSV writer & \Done &
  Removed the 25--220\,ms periodic stalls. \\
ABBA ordering & \Open & Implemented, never tested on hardware. \\
\bottomrule
\end{tabular}
\end{center}

\section{Round 2 --- 2026-08-16/17, against the 2026-08-14 session}

\begin{center}
\small
\begin{tabular}{@{}llp{8.2cm}@{}}
\toprule
\hd{Commit} & \hd{Scope} & \hd{What it did} \\
\midrule
\co{a98fd2f} & \co{data:} &
  2026-08-14 session (101\,MB): four averaged runs, one stream run, one burst run
  all at $\tau$=\SI{120}{\micro\second}; a 12-point integration-time sweep
  (1--\SI{213}{\micro\second}, 5 repeats each); an 8-peak reference sweep; a
  multipoint run. Committed \emph{before} any code change, so the analysis has a
  fixed baseline. \\[3pt]
\co{6eb7eb5} & \co{docs:} &
  \textbf{Retraction} of the ``rate is FPGA-bound'' conclusion in the 2026-08-06
  document, with the evidence, plus marking \S4.3's root cause disproved. \\[3pt]
\co{524dbb4} & \co{feat:} &
  \file{twopoint_spectra.py} (365 lines) and \file{twopoint_postprocess.py}
  (790 lines); \co{clean_burst_lockin.py} rewritten to delegate. Found and fixed the
  per-channel despiking bug and the 12\% burst-cadence error. \\[3pt]
\co{6893bcb} & \co{perf:} &
  Corrected rate model (\co{HOST_CALL_S}/\co{HOST_FLOOR_S} replacing
  \co{HOST_OVERHEAD_S}); \co{reps_that_fit()}; \co{measure_host_floor()};
  \co{multipoint_on_stale="drop"}; \co{stream()} sequence aligned with
  \co{acquire()} and packet-size validation. \\[3pt]
\co{8fd303f} & \co{refactor:} &
  Notebook split into Step 4A/5A, 4B/5B, 4C/5C, Step 6. New
  \file{twopoint_runner.py} (236 lines) for the shared plumbing. 1497 lines of
  notebook changed. \\[3pt]
\co{68e95a9} & \co{feat:} &
  Three rig tests: \co{--compare-read-paths}, \co{--stream-scan}, \co{--floor}. \\[3pt]
\co{f1bd0eb} & \co{docs:} &
  The 2026-08-14 methods document, \co{analyze_twopoint_0814.py}, five generated
  tables and three figures; pointer fixes in \co{LOCKIN_ACQUIRE.md} and
  \co{twopoint_lockin_notes.md}. \\
\bottomrule
\end{tabular}
\end{center}

\section{Round 3 --- 2026-08-17, the board wedge}

Triggered by a live failure at the rig rather than by analysis: Step~4C hung at its
auto-zero \co{acquire()}, and \co{Ctrl-C} would not release it. Diagnosed to
\co{poll\_data()}'s \co{timeout=None} default and fixed the same day. Full account in
\S\ref{sec:wedge}; the naming collision in \S\ref{sec:namecollision} was found while
tracing which file each mode had written.

\begin{retractbox}[title=\textbf{Read Round 4 before acting on anything in this section}]
The central change below --- bounding every poll --- was \textbf{withdrawn the same
day}. It broke burst and streaming on first contact with the board.
\S\ref{sec:boundedpollretraction} is the account; \S\ref{sec:round4} is what
replaced it. The interrupt cleanup and the naming split survive unchanged.
\end{retractbox}

\begin{center}
\small
\begin{tabular}{@{}lp{10.4cm}@{}}
\toprule
\hd{Scope} & \hd{What it did} \\
\midrule
\co{fix(qickdawg):} &
  \co{acquire()}'s drain loop extracted to \co{\_drain\_readout()} and bounded: an
  explicit \co{poll\_data} timeout, a no-progress deadline that raises
  \co{TimeoutError}, and a \co{BaseException} handler that stops the tProc and the
  readout before re-raising. Adds \co{POLL\_TOTALTIME\_S}, \co{POLL\_TIMEOUT\_S},
  \co{ACQUIRE\_IDLE\_S}, \co{\_acquire\_timeout\_s()} and \co{\_abort\_readout()}. \\[3pt]
\co{fix(twopoint):} &
  \co{stream()} passes the same explicit timeout --- which is what finally makes its
  own \co{poll\_timeout\_s} reachable --- and wraps its loop in \co{try/finally}, so an
  interrupted or abandoned generator leaves the board idle. Records
  \co{stream\_qc()["aborted"]}. \\[3pt]
\co{feat(twopoint):} &
  \co{MultipointLockinODMR.recover()} and \co{ensure\_idle()}. The pre-flight is wired
  into Steps 4A, 4B and 4C. \\[3pt]
\co{fix(notebook):} &
  Step~4A writes \co{twopoint\_lockin\_avg\_*} and Step~4B \co{twopoint\_lockin\_burst\_*}
  instead of both writing \co{\_live\_}; \co{find\_live\_csvs()} accepts all three
  prefixes. \\
\bottomrule
\end{tabular}
\end{center}

\subsection{Verified offline}

No board was available, so all three paths were exercised against fakes built from
the real \co{qick} sources:

\begin{center}
\small
\begin{tabular}{@{}>{\raggedright\arraybackslash}p{2.4cm}
                  >{\raggedright\arraybackslash}p{6.4cm}
                  >{\raggedright\arraybackslash}p{4.6cm}@{}}
\toprule
\hd{Check} & \hd{What it drives} & \hd{Result} \\
\midrule
Root cause & real \co{DataStreamer}, counter frozen;
  \co{poll\_data()} vs \co{poll\_data(timeout=1.0)} &
  hangs vs returns in 1.01\,s \\[3pt]
Healthy drain & \co{\_drain\_readout()} from the edited file, 50 shots in
  10-shot packets & 50/50, buffer bit-exact \\[3pt]
Stalled tProc & same, board returns nothing &
  \co{TimeoutError} at the budget, board cleaned \\[3pt]
Interrupt & \co{KeyboardInterrupt} mid-drain &
  re-raised, board cleaned \\[3pt]
Stream abort & real \co{stream()} source, consumer breaks after 10 reps &
  \co{finally} ran, \co{aborted=True} \\[3pt]
Stream completion & same, run to 20/20 reps &
  cleaned, \co{aborted=False} \\
\bottomrule
\end{tabular}
\end{center}


\section{Round 4 --- 2026-08-17 evening, undoing Round 3's central change}
\label{sec:round4}

Round 3 shipped at 14:23 and the rig session at 14:25--14:29 showed it had broken
two of the three acquisition methods. The evidence is the notebook's own saved
outputs:

\begin{center}
\small
\begin{tabular}{@{}cllll@{}}
\toprule
\hd{Cell} & \hd{Step} & \hd{exec} & \hd{reps in the failing call} & \hd{Result} \\
\midrule
18 & 4A averaged & 21 & 35   & worked --- 1067 samples, \SI{104.5}{\hertz} \\
22 & 4B burst    & 25 & 2000 & \co{TimeoutError}: read 0 of 2000 shots in \SI{10}{\second} \\
26 & 4C stream   & 27 & 2000 & hung in \co{Pyro4 socketutil.recvall}; released by \co{Ctrl-C} \\
24 & 5B analysis & 47 & ---  & \co{TypeError: ... not 'NoneType'} (collateral) \\
\bottomrule
\end{tabular}
\end{center}

Both real failures are the dedicated large-reps auto-zero \co{acquire()}, which
4A never makes. The comparison that settled it was the 2026-08-17 \emph{morning}
session, taken on a second computer at \co{8fd303f} --- before Round 3 --- which
streamed \num{31250} samples at \SI{1041.7}{\hertz} with zero missing reps, four
times over. Same rig, same board, same day; the only difference was the code.

\begin{center}
\small
\begin{tabular}{@{}lp{10.4cm}@{}}
\toprule
\hd{Scope} & \hd{What it did} \\
\midrule
\co{revert(qickdawg):} &
  \co{POLL\_TIMEOUT\_S} back to \co{None}, restoring qick's blocking
  \co{poll\_data()}. The no-progress deadline is now consulted only when the poll
  is bounded, which it no longer is by default; set the class attribute to a float
  for a diagnostic run. \S\ref{sec:boundedpollretraction} carries the reason at
  length so this is not re-attempted. \\[3pt]
\co{fix(qickdawg):} &
  An empty return under a blocking poll now raises, naming how many shots were
  read. Previously that condition spun on the board with RPCs; the naive
  alternative would short-read zeros into the tail of \co{d\_buf}. \\[3pt]
\co{fix(twopoint):} &
  \co{ensure\_idle()} reports and no longer calls \co{recover()} implicitly.
  \co{recover()} is staged: stop, drain the data queue, and only escalate to
  \co{start\_worker()} if the board is still not idle (\co{force=True} skips
  ahead). Steps 4A/4B/4C raise and point at the new Step~0 instead of continuing
  into a hang. \\[3pt]
\co{fix(twopoint):} &
  \co{stream\_headroom()} read \co{qd.soc['readouts']}, and \co{qd.soc} is a Pyro
  \co{Proxy} --- not subscriptable --- so the check printed an error for the whole
  2026-08-17 session instead of a number. It now uses \co{qd.soccfg}. This is the
  check the multipoint streaming path actually depends on. \\[3pt]
\co{fix(notebook):} &
  \co{twopoint\_postprocess.process(None)} raises a named error saying which run
  cell to re-run; Steps 5A/5B/5C fall back to the newest matching CSV on disk.
  \co{default\_config.mw\_gain} hoisted to a named \co{MW\_GAIN} constant with its
  provenance --- it went \num{10500}$\rightarrow$\num{15000} inside a commit whose
  subject was ``normalise JSON encoding''. \\[3pt]
\co{feat(notebook):} &
  Steps 4S and 5S in \file{Modules/Lockin_module.ipynb}: streaming for the
  multipoint lock-in with its own post-processing and vector reconstruction
  (Ch.~\ref{ch:mpstream}). \\[3pt]
\co{test(scripts):} &
  \file{scripts/test\_drain\_readout.py}, which exercises the drain loop against
  \co{qick}'s \emph{real} \co{poll\_data} rather than a mock --- the gap that let
  Round 3 ship. 17 checks, no board needed. \\[3pt]
\co{docs(scripts):} &
  \file{scripts/analyze\_twopoint\_0817.py} regenerates every 2026-08-17 number in
  this document. \\
\bottomrule
\end{tabular}
\end{center}

\begin{keybox}[The process lesson, which is the expensive one]
Round 3's commit message says ``Verified offline (no board)''. It was --- against a
mock written alongside the fix, which modelled \co{poll\_data} as ``returns
packets, or \co{[]} on timeout''. The real \co{poll\_data} also discards a packet
when \co{stop\_flag} is set, and returns \co{[]} the moment its \SI{100}{\milli\second}
\co{totaltime} expires. Neither behaviour was in the mock, and both matter.

\file{scripts/test\_drain\_readout.py} now extracts \co{poll\_data}'s source from
the installed \co{qick} and executes it, and fails loudly if that extraction stops
matching. A mock of the component you are debugging tests your belief about it,
not it.
\end{keybox}

\begin{openbox}[title=\textbf{Two more traps found while tracing this, both still live}]
\begin{itemize}
  \item \textbf{\co{qickdawg} resolves to the wrong checkout unless the project
        root is first on \co{sys.path}.} It is pip-installed editable from
        \file{\textasciitilde/Documents/PhD/NV Compact Magnetometer/qick-dawg/src},
        which carries none of this project's patches. The notebooks are safe --- cell~0
        inserts \co{PROJECT\_ROOT} at position 0 --- but a bare
        \co{python scripts/foo.py} is not. \file{test\_drain\_readout.py} asserts
        which copy it imported; other scripts do not.
  \item \textbf{The burst replay reproduces on a second machine.} 50.0--50.8\%
        across six 2026-08-17 morning burst files, at \co{8fd303f}. It is not an
        artefact of this laptop, and it is not the wedge (\S\ref{sec:replay}).
\end{itemize}
\end{openbox}

\section{The notebook layout}
\label{sec:notebook}

\file{Modules/Twopoint_Lockin_module.ipynb}, 34 cells. Step 4 used to be a single
26\,kB cell with a \co{LIVE_BURST_MODE} switch, and Step 5 a single generic analysis
cell that treated all three methods alike --- which is wrong in a different way for
each of them.

\begin{center}
\small
\begin{tabular}{@{}clp{8.4cm}@{}}
\toprule
\hd{Cells} & \hd{Section} & \hd{Contents} \\
\midrule
0--5   & setup & project root, \co{MW_GAIN} and \co{default_config} ($\tau$ set
  here), DSA attenuation \\
6--7   & \textbf{Step 0} & board recovery. Skip it on a normal day; run it after
  any interrupted acquisition (\S\ref{sec:wedge}) \\
8--10  & ODMR sweep & the reference sweep \\
11--14 & Step 1a / 1b & pick the peak to track; place the parked pair and build
  \co{TWOPOINT_CALIB} (this is where the slopes and baselines come from) \\
15--18 & Step 2 / 3 & one batch at the two parked frequencies; apply the conversion \\
\midrule
19--20 & \textbf{Step 4A} & averaged acquisition. Measures the host floor, sizes
  \co{reps} to fill it, prints the alias warning. Writes \co{twopoint_lockin_avg_*} \\
21--22 & Step 5A & \co{process_averaged} + ASD, $\sim$15 lines \\
23--24 & \textbf{Step 4B} & burst acquisition, \co{on_stale="drop"}, cadence-based
  timestamps, reports \co{BURST_STALE_FRACTION}. Writes \co{twopoint_lockin_burst_*} \\
25--26 & Step 5B & \co{process_burst} + per-segment ASD \\
27--28 & \textbf{Step 4C} & streaming acquisition, \co{stream_qc()} assertions.
  Writes \co{twopoint_lockin_stream_*} \\
29--30 & Step 5C & \co{process_stream} + ASD, optional mains notch \\
31--32 & Step 6 & post-process \emph{any} saved run; mode detected from the columns;
  side-by-side comparison table \\
33     & Caveats & what the method does and does not measure \\
\bottomrule
\end{tabular}
\end{center}

Each Step 4x opens with a one-RPC \co{ensure_idle()} that raises and points at
Step~0 rather than starting into a board that is not free. Each Step 5x falls back
to the newest run of its own mode on disk when its partner's \co{*_LAST_CSV} is
unset, so a failed acquisition produces a sentence rather than a \co{TypeError}
from inside \co{pathlib}.

Each Step 4x is self-contained --- it builds its own program, primes, auto-zeroes and
writes its own CSV --- with the parameters that change run to run inline, where they
are edited. Each Step 5x reads the run its partner just wrote. The shared mechanical
plumbing (append-only CSV writer, row building, the timing/freshness/calibration
reports, the standard figure) lives in \file{notebook_modules/twopoint_runner.py} so
that 4A and 4B do not carry 150 duplicated lines each.

\section{Round 6 --- 2026-08-18, the shot counter was never cleared}
\label{sec:counter}

\begin{keybox}[One missing line, and it was never in this file]
\co{acquire()} and \co{stream()} arm the board-side readout without first zeroing
the tProc shot counter. qick's own \co{AcquireProgram.start_round()} does it one
line after \co{reload_mem()}:
\begin{quote}\small\ttfamily
soc.reload\_mem()\\
\# make sure count variable is reset to 0\\
soc.clear\_tproc\_counter(addr=self.counter\_addr)
\end{quote}
The vendored \file{qickdawg/nvpulsing/nvaverageprogram.py} reproduces every other
step of that sequence and drops this one.
\end{keybox}

\subsection{What the counter does}

It is tProc \emph{data memory} at \co{counter_addr}, and nothing in the acquire
path resets it: the program only ever increments it, \co{stop_tproc(lazy=True)}
is a documented no-op on tProc v1 (\S\ref{sec:rpc}, stage 1b), and
\co{reload_mem()} reloads the program's own data, not this address. So it carries
the previous run's final value into the next run.

The board-side worker (\file{qick/streamer.py}) then does:

\begin{quote}\small\ttfamily
shots = self.soc.get\_tproc\_counter(addr=counter\_addr)\\
if shots >= min(last\_shots+stride, total\_shots):\\
\ \ \ \ newshots = shots - last\_shots\quad\#\ last\_shots starts at 0\\
\ \ \ \ data = self.soc.get\_accumulated(ch=ch, address=addr,\\
\ \ \ \ \ \ \ \ \ \ \ \ \ \ \ \ \ \ \ \ \ \ \ \ \ \ \ \ \ \ \ \ \ \ \ \ length=newshots*reads\_per\_count)
\end{quote}

With a stale counter that test passes on the worker's \emph{first} look, before
the new program has written anything, and it transfers whatever is sitting in the
accumulated buffer. Two outcomes, depending only on how the stale value compares
to the new run's length:

\begin{center}
\small
\begin{tabular}{@{}lp{9.6cm}@{}}
\toprule
\hd{Stale counter} & \hd{What happens} \\
\midrule
$>$ \co{total_count} &
  One oversized packet. \textbf{This is the 2026-08-18 crash}: a 100-rep priming
  \co{acquire()} inherited 2000 from the preceding auto-zero and raised
  \co{ValueError: could not broadcast input array from shape (4000,2) into shape
  (200,2)} --- $2000\times2$ readouts into a $100\times2$ array. \\[3pt]
$=$ \co{total_count} &
  \textbf{No error at all.} The call returns in milliseconds with the
  \emph{previous} run's contents. \\[3pt]
$<$ \co{total_count} &
  A short head of stale data, then the run proceeds normally. \\
\bottomrule
\end{tabular}
\end{center}

\subsection{Which paths were affected}

\co{reset_tproc=True} routes \co{config_all} through \co{stop_tproc(lazy=False)},
which on tProc v1 is \co{tproc.reset()} + \co{reload_program()} --- and a reset
plausibly zeroes data memory as a side effect. That is why the damage is uneven:

\begin{center}
\small
\begin{tabular}{@{}lcp{6.4cm}@{}}
\toprule
\hd{Call site} & \hd{\co{reset_tproc}} & \hd{Exposure} \\
\midrule
Step 4A, every sample            & default \co{False} & every call after the first inherited the last one's count \\
Step 4B, the burst \co{acquire}  & \textbf{\co{True}} & probably protected by the reset \\
Step 4B/4C priming and auto-zero & default \co{False} & \textbf{where the crash happened} \\
Step 4C \co{stream()}            & n/a                & armed \co{start_readout} directly, never cleared \\
ODMR sweep, Steps 1--3           & default \co{False} & one \co{acquire} per sweep point \\
\bottomrule
\end{tabular}
\end{center}

\begin{openbox}[Do not yet credit this with the burst replay]
It is the obvious suspect for \S\ref{sec:replay} --- ``the call returns in
\SI{13}{\milli\second} instead of \SI{425}{\milli\second} with 95.7\% of its rows
bit-identical to the batch before'' is exactly the middle row of the table above.
But burst mode is the one path that passes \co{reset_tproc=True}, so if the v1
reset does zero the counter, burst was already protected and its replay has
another cause.

The rig settles it in one line, run straight after a burst batch:
\begin{center}\co{qd.soc.get_tproc_counter(addr=1)}\end{center}
Zero means the reset clears it and the replay is something else. Non-zero means
this was the cause all along, and \S\ref{sec:replay} can be closed. Step 0's
\co{recover()} now prints the value for exactly this reason.
\end{openbox}

\subsection{The fix}

\begin{center}
\small
\begin{tabular}{@{}lp{8.8cm}@{}}
\toprule
\hd{Where} & \hd{Change} \\
\midrule
\co{acquire()} & \co{clear_tproc_counter(addr=self.counter_addr)} between
  \co{reload_mem()} and \co{start_readout()}, matching qick. One RPC,
  $\sim$\SI{1}{\milli\second}, on a call that already costs 10--11\,ms. \\[3pt]
\co{stream()} & the same, in the same position. \\[3pt]
\co{_drain_readout()} & an oversized packet now raises a named \co{RuntimeError}
  that says the counter did not start at zero, instead of letting numpy raise
  \co{could not broadcast}, which names neither cause nor fix. \\[3pt]
\co{_abort_readout()} & clears it on the way out, so an interrupted cell does not
  hand a poisoned counter to whatever runs next. \\[3pt]
\co{recover()} & reads the counter, prints it, clears it. \\
\bottomrule
\end{tabular}
\end{center}

\textbf{No notebook cell changed.} Every acquisition in both notebooks goes
through \co{acquire()} or \co{stream()}; the vendored qickdawg has exactly one
\co{start_readout} call site and \co{stream()} is the other, and every
\co{LockinODMR}-style subclass delegates to \co{NVAveragerProgram.acquire} via
\co{super()}.

\subsection{Verified}

\file{scripts/test_shot_counter.py}, 21 checks. It drives qick's \emph{real}
\co{DataStreamer._run_readout} against a fake board whose counter persists between
runs and whose program ramps the counter rather than completing instantly --- the
arithmetic under test lives in that worker, so a mock of the worker would prove
nothing. It reproduces the exact reported shapes, \co{(4000,2)} into
\co{(200,2)}, and shows the clear removes them; shows the equal-counter case
completing silently with data read before the program produced it; and confirms
the call is present in the shipped \co{acquire()}, \co{stream()},
\co{_abort_readout()} and \co{recover()}, in the right order relative to
\co{reload_mem} and \co{start_readout}.

\section{Round 7 --- 2026-08-18, revert the acquisition path to \texorpdfstring{\co{8fd303f}}{8fd303f}}
\label{sec:revert}

Three rounds of changes to the acquisition path each fixed something real and each
left the rig worse than it found it. On 2026-08-18 the operator reported that
streaming still did not run and that averaged mode had fallen to \SI{48}{\hertz},
against the 87--105\,Hz it had held all week, and asked for the path to be made
identical to the revision on the second machine.

\begin{center}
\small
\begin{tabular}{@{}llp{6.0cm}@{}}
\toprule
\hd{File} & \hd{State} & \hd{Why} \\
\midrule
\file{qickdawg/nvpulsing/nvaverageprogram.py} & \textbf{byte-identical to \co{8fd303f}} & the acquire loop and the poll \\
\file{notebook_modules/multipoint_lockin_program.py} & \textbf{byte-identical to \co{8fd303f}} & \co{stream()}, the timing model \\
\emph{all 35} \file{qickdawg/*.py} & \textbf{identical} & nothing else on the path drifted \\
\midrule
\file{Modules/*.ipynb} & \textbf{kept} & the revert is the acquisition \emph{modules}; the notebooks keep Step 0, the \co{_avg_}/\co{_burst_}/\co{_stream_} names, \co{mw_gain = 15000}, the newest-run fallbacks, Steps 4S/5S and figure saving \\
\file{notebook_modules/twopoint_postprocess.py} & kept & post-processing only; never on the acquire path \\
\bottomrule
\end{tabular}
\end{center}

The only notebook edit the revert forced: Step 0 and the three pre-flight guards
called \co{recover()} and \co{ensure_idle()}, which live in the acquisition module
and went away with it. Both are now written inline against the raw \co{QickSoc}
proxy --- \co{stop_tproc}, \co{streamer.stop_readout}, \co{poll_data(totaltime=-1)},
\co{streamer.readout_running} --- all plain remote methods qick itself exposes. So
the notebook keeps the escape hatch and is independent of which module revision is
loaded. The guard also downgraded from fatal to a warning on probe failure: a
board-state probe must never be able to stop a run.

\subsection{What went away with it}

\begin{center}
\small
\begin{tabular}{@{}lp{8.6cm}@{}}
\toprule
\hd{Gone from the modules} & \hd{Consequence} \\
\midrule
\co{POLL_TIMEOUT_S}, \co{_drain_readout} & the poll is qick's plain
  \co{poll_data()} again --- blocking, which is the behaviour that worked. \\
\co{_abort_readout()} & an interrupted cell no longer stops the tProc on the way
  out. After a Ctrl-C mid-readout, run Step 0; if it still reports busy, restart
  the board's Pyro server. A kernel restart will not clear it. \\
\co{ensure_idle()}, \co{recover()} & reimplemented in the notebook instead, as
  above. \\
\co{clear_tproc_counter()} & the stale-counter behaviour of \S\ref{sec:counter} is
  live again --- see the box below. \\
\bottomrule
\end{tabular}
\end{center}

The hardware operating point is \emph{unchanged}: \co{mw_gain = 15000},
$\tau = \SI{120}{\micro\second}$, \co{AVG_REPS_PER_BATCH = None},
\co{BURST_REPS = 500}, \co{STREAM_TARGET_HZ = 1000}. Those live in the notebook,
not in the reverted modules.

\begin{openbox}[The stale shot counter is real and is now unfixed]
\S\ref{sec:counter} is not retracted --- the mechanism was reproduced against
qick's own \co{DataStreamer._run_readout}, and \co{8fd303f} is missing the same
call qick makes. It is an open defect of the reverted revision, and it is the
leading candidate for the burst replay of \S\ref{sec:replay}.

\file{scripts/test_shot_counter.py} is kept for that reason, scoped down to the
mechanism: 13 checks, all passing, none of them asserting anything about the
shipped code. If the fix is revisited, that test is the check it needs to pass.
\end{openbox}

\subsection{The lesson this round is actually about}

Every one of Rounds 3, 4 and 6 was reasoned from the source, verified offline
against a test, and shipped --- and each one changed the rate or broke a mode on
first contact with the rig. Round 3 was verified against a mock and broke burst
and streaming. Round 6 was verified against qick's real worker thread and, whatever
else it did, coincided with averaged mode falling to \SI{48}{\hertz}.

Offline verification says the code does what you think. It says nothing about
whether what you think is worth doing. \textbf{On this rig, the acquisition path
changes one thing at a time, with a before-and-after run on the bench, or it does
not change.}

\begin{center}
\small
\begin{tabular}{@{}lrl@{}}
\toprule
\hd{Session} & \hd{Averaged rate} & \hd{Path revision} \\
\midrule
2026-08-14 & 86.8--88.8\,Hz & \co{8fd303f} lineage \\
2026-08-17 morning \emph{(other machine)} & --- & \co{8fd303f} \\
2026-08-17 afternoon & 104.5\,Hz & \co{9832e07} (burst + stream broken) \\
2026-08-18 & \textbf{48\,Hz} & Round 6 \\
2026-08-18, after this revert & \emph{to be measured} & \co{8fd303f} \\
\bottomrule
\end{tabular}
\end{center}

The last row is the first thing to fill in at the rig.

\section{Round 5 --- 2026-08-18, every figure goes to disk}
\label{sec:autosave}

Neither notebook wrote a single PNG. Every plot lived only in cell output, which
"Restart \& Clear Outputs" discards and which cannot be opened next to the CSV it
describes. The \co{*_result.png}/\co{*_fft.png} pairs already sitting in
\file{data/results/} were produced by a separate tool on the acquisition machine,
not by these notebooks --- nothing in this repository wrote them.

\subsection{\texorpdfstring{\co{notebook_modules/figure_autosave.py}}{figure\_autosave.py}}

Saving is automatic rather than a \co{savefig} line per plot, because there are
twenty-three plotting sites across the two notebooks and adding a line to each
guarantees the twenty-fourth is forgotten.

The module hooks the two moments at which the inline backend destroys a figure:

\begin{center}
\small
\begin{tabular}{@{}lp{9.2cm}@{}}
\toprule
\hd{Hook} & \hd{What it catches} \\
\midrule
\co{post_execute} callback &
  Figures a cell drew but never explicitly showed. It must run \emph{before}
  matplotlib-inline's \co{flush_figures}, which displays and then closes
  everything; IPython fires \co{post_execute} in registration order, so
  \co{enable()} registers its own callback and then re-registers
  \co{flush_figures} to move it to the back of the queue. \\[3pt]
wrapper around \co{pyplot.show} &
  Figures closed mid-cell. \co{show()} on this backend closes what it displays,
  and \co{twopoint_runner.plot_run}, \co{TwoPointResult.figure(show=True)} and the
  live loops in \file{Lockin_module.ipynb} all call it. \\
\bottomrule
\end{tabular}
\end{center}

Both funnel into \co{capture()}, and the ``already written'' marker lives on the
figure object, so the two hooks cannot duplicate each other's work.

\subsection{Naming}

\co{set_run(csv, tag=...)} points the module at a run: the PNGs take the CSV's
stem and land in the CSV's own folder. The numbering restarts on every
\co{set_run}, so \textbf{re-running a cell overwrites that cell's own PNGs}
instead of accumulating copies.

\begin{center}
\small
\begin{tabular}{@{}ll@{}}
\toprule
\hd{Cell} & \hd{Writes} \\
\midrule
Step 4A/4B/4C \co{set_run(AVG_CSV, tag="run")} & \co{twopoint_lockin_avg_<ts>_run_result.png} \\
Step 5A/5B/5C \co{set_run(A5_CSV)} & \co{twopoint_lockin_avg_<ts>_result.png} \\
 & \co{twopoint_lockin_avg_<ts>_fft.png} \\
\bottomrule
\end{tabular}
\end{center}

A figure is classified \co{_fft} only when \emph{all} of its axes are log--log
spectra. That matters here: the three-panel figure from
\co{TwoPointResult.figure()} carries an ASD as its last panel, and must not claim
the \co{_fft} name from the standalone spectrum figure of the same run. A cell can
override the classification entirely with \co{fig.set_label("droop")}, giving
\co{<stem>_droop.png}.

\subsection{\texorpdfstring{\co{TwoPointResult.spectrum_figure()}}{spectrum\_figure()}}

Added so that every two-point run gets a real \co{_fft.png} rather than only the
ASD panel buried at a third height inside \co{figure()}. Full width, with the
white floor, $\eta$ from $\sigma$, and every line above 2$\times$ the floor
annotated with its excess. Steps 5A/5B/5C call it after \co{figure()}.

\subsection{Verified}

\begin{itemize}
  \item \file{scripts/test_figure_autosave.py} --- 18 checks, driven through a
        real \co{InteractiveShell} with \co{\%matplotlib inline} rather than a
        mock, because hook ordering against \co{flush_figures} is the whole
        mechanism. Confirms the ordering explicitly, and covers the unshown-figure
        path, the \co{plt.show()} path, the mixed-figure classification, tags,
        overwrite-on-rerun, empty figures, \co{disable()}/re-\co{enable()} and the
        label override.
  \item \file{scripts/test_notebook_figure_saving.py} --- runs Step 5B and Step 5C
        \emph{out of the \co{.ipynb} files} against scratch copies of the
        2026-08-17 runs, and checks both PNGs appear beside the CSV and that a
        second run overwrites them. 6 checks.
  \item A static pass over both notebooks asserts every cell that draws calls
        \co{set_run} \emph{before} its first plot. That caught one cell where the
        call had landed after the plots and would have filed them under the
        previous run's name.
\end{itemize}

\begin{openbox}[The live loops write one figure, not a movie]
\file{Lockin_module.ipynb} Steps 4 and 4B redraw a single figure in place while
the run proceeds. What reaches disk is that figure's \emph{final} state, which is
the whole run --- not a frame per update. Step 4BN and Step 4S are headless and
draw once at the end, so they are unaffected.

There is also no \co{_fft.png} for the multipoint live cells: their time base is
the host loop period, not a cadence, so an ASD computed on it would put a
frequency scale on jitter. The streamed Step 5S has a real cadence and does
produce one.
\end{openbox}

\section{The code, file by file}
\label{sec:filemap}

\begin{center}
\small
\begin{longtable}{@{}p{5.0cm}p{2.0cm}p{7.2cm}@{}}
\toprule
\hd{File} & \hd{State} & \hd{Role} \\
\midrule
\endhead
\multicolumn{3}{@{}l}{\emph{Acquisition}} \\
\file{notebook_modules/} \file{multipoint_lockin_program.py} & modified &
  The FPGA program, all three read paths, the timing model, the freshness guard. \\
\file{qickdawg/nvpulsing/} \file{nvaverageprogram.py} & modified &
  Vendored; the \co{config_all} fix and \co{reset_tproc} plumbing. \\
\file{Modules/} \file{Twopoint_Lockin_module.ipynb} & restructured &
  32 cells: setup, ODMR, Steps 1--3, then 4A/5A, 4B/5B, 4C/5C, 6. \\
\midrule
\multicolumn{3}{@{}l}{\emph{Analysis (new)}} \\
\file{notebook_modules/} \file{twopoint_spectra.py} & \textbf{new} &
  Every FFT. One-sided PSD, $\eta=\sigma\sqrt{2\Delta t}$, mains-line finder,
  alias report, comb notch. Verified against Parseval. \\
\file{notebook_modules/} \file{twopoint_postprocess.py} & \textbf{new} &
  One pipeline per mode + mode detection + exact readout-quantum recovery + CLI. \\
\file{notebook_modules/} \file{twopoint_runner.py} & \textbf{new} &
  Shared acquisition plumbing: append-only CSV, row building, the three reports,
  the standard figure. \\
\file{notebook_modules/burst_qc.py} & reused &
  Batch table, stale masks, retiming, segment PSD, block averaging. \\
\file{notebook_modules/} \file{spike_rejection.py} & reused &
  \co{HampelDespiker}. \\
\midrule
\multicolumn{3}{@{}l}{\emph{Scripts}} \\
\file{scripts/} \file{analyze_twopoint_0814.py} & \textbf{new} &
  Regenerates every number and figure in this document. \\
\file{scripts/} \file{profile_twopoint_acquire.py} & modified &
  Static signature check + per-RPC timing + the three rig tests. \\
\file{scripts/} \file{clean_burst_lockin.py} & rewritten &
  Burst CLI, now delegating to \co{twopoint_postprocess}. \\
\file{scripts/} \file{analyze_twopoint_session.py} & unchanged &
  The 2026-08-06 analysis; still regenerates that document. \\
\midrule
\multicolumn{3}{@{}l}{\emph{Documents}} \\
\file{docs/twopoint_master_reference/} & \textbf{new} & This document. \\
\file{docs/} \file{2026-08-14_twopoint_methods/} & new & Markdown source + generated tables/figures. \\
\file{docs/} \file{2026-08-06_twopoint_timing/} & annotated & The earlier analysis, with corrections marked in place. \\
\file{docs/LOCKIN_ACQUIRE.md} & pointer fixed & Repeated the retracted 1.6\,ms figure. \\
\file{docs/} \file{twopoint_lockin_notes.md} & pointer fixed & Same. \\
\bottomrule
\end{longtable}
\end{center}

\section{Wiki}

The synthesis is filed in the Obsidian vault at \file{NV Project/}.

\begin{itemize}
  \item \textbf{New page:} \file{wiki/sources/twopoint_three_modes_2026_08_14.md}
  \item \textbf{Correction banner:} \file{wiki/sources/twopoint_timing_and_burst_defect.md}
  \item \textbf{Updated:} \file{wiki/experiments/lockin_tracking.md} (new
        ``August 14--17'' section, and the August 6--12 one partly retracted);
        \file{wiki/concepts/sensitivity.md} (three new rows, plus a warning that the
        two-point noise does not average down);
        \file{wiki/questions/open_questions.md} (three questions answered, four new
        ones with pass criteria); \file{wiki/overview.md} (status table and the
        ``Where We Stand'' list); \file{index.md} (routing).
  \item \textbf{Log entry:} \file{log.md}, dated 2026-08-17.
\end{itemize}

\emph{Not} committed to git: the vault sits inside the home-directory repository,
which contains unrelated and sensitive untracked files. The files are on disk and
sync through OneDrive.


\chapter{What still needs to be done}
\label{ch:todo}

\section{Rig checks, in priority order}
\label{sec:rigchecks}

Everything here needs the RFSoC at 192.168.0.103, reachable only from the Windows
rig machine. Each has a pass criterion.

\begin{center}
\small
\begin{longtable}{@{}p{0.7cm}p{4.6cm}p{4.4cm}p{4.6cm}@{}}
\toprule
\hd{\#} & \hd{Check} & \hd{Command} & \hd{Pass criterion} \\
\midrule
\endhead
S0 & \textbf{Does the wedge fix close the burst replay?} Now the cheapest
  decisive test, because \S\ref{sec:wedge} gives \S\ref{sec:replay} its first
  mechanism since \co{reset_tproc} was eliminated. \emph{Do this first.}
  $\sim$1\,min. &
  Run Step~4B for 30\,s and read \co{BURST_STALE_FRACTION}. &
  \textbf{0\%} $\Rightarrow$ the replay was the wedge and both defects close.
  Still $\sim$\textbf{50\%} $\Rightarrow$ they are independent; S1 next. \\[3pt]
S1 & \textbf{Burst vs stream, interleaved.} Decides whether the $\sim$6$\times$
  streaming excess is the read path or the room. \emph{Do this first.} $\sim$1\,min. &
  \co{profile_twopoint_acquire.py --compare-read-paths --tau 120 --reps 500} &
  ASD ratio within \textbf{1.25$\times$} $\Rightarrow$ environment. Well above
  $\Rightarrow$ software defect. \\[3pt]
S2 & Does the streaming excess grow with run length (buffer fill)? &
  \co{profile_twopoint_acquire.py --stream-scan} &
  Floor at 125\,k reps within \textbf{1.3$\times$} of the floor at 2\,k reps. \\[3pt]
S3 & Confirm the per-call floor and the flat rate. &
  \co{profile_twopoint_acquire.py --floor --tau 120 --reps 1 4 10 23 42 100} &
  Fastest call \textbf{9--11.5\,ms} with a \co{reps=1} probe; period flat to
  \textbf{15\%} across reps whose FPGA work fits under it. \\[3pt]
S4 & \textbf{Fill in \S\ref{sec:rpc}} --- the per-RPC split, the one structural gap
  in this document. &
  \co{profile_twopoint_acquire.py --ip 192.168.0.103} &
  A time for each of stages 1--7. No pass/fail; this is a measurement that has never
  been made. \\[3pt]
S5 & Free averaging. 30\,s Step 4A with \co{AVG_REPS_PER_BATCH = None}. &
  Notebook Step 4A &
  Still $\sim$88\,Hz \emph{and} $\sigma(\Delta f)\le$ \textbf{54.6\,kHz}. Given
  \S\ref{sec:notshot}, expect less than the naive 2$\times$. \\[3pt]
S6 & \textbf{The per-acquire noise term.} Interleave averaged runs at reps =
  10 / 23 / 42 within one sitting. &
  Notebook Step 4A, three configs, alternating &
  Does $\sigma$ fall as $1/\sqrt{\text{reps}}$? If not, the term is per-call --- then
  test whether it tracks the idle gap by varying \co{SAVE_EVERY_SEC} or inserting a
  deliberate sleep. \\[3pt]
S7 & ABBA A/B at fixed $\tau$ and equal averaging (\co{"forward"} 46 reps vs
  \co{"abba"} 23 reps), interleaved. &
  Notebook Step 4A, two configs &
  Report the change in $\sigma(\Delta f)$. Any reduction is free. \\[3pt]
S8 & Honest burst rate with the drop policy. 10\,s Step 4B. &
  Notebook Step 4B, \co{BURST_ON_STALE="drop"} &
  Step 5B reports \textbf{\co{n_stale_batches} = 0} on the file Step 4B just wrote,
  and \co{BURST_STALE_FRACTION} is printed. \\[3pt]
S9 & Instrument the replay. Print the shot counter and queue depth before/after each
  \co{start_readout} during a burst run. &
  needs a small patch to \co{acquire()} &
  Discriminates the three candidate mechanisms in \S\ref{sec:replay}. \\
\bottomrule
\end{longtable}
\end{center}

\section{Software backlog}

\begin{center}
\small
\begin{tabular}{@{}p{5.4cm}lp{7.2cm}@{}}
\toprule
\hd{Item} & \hd{Priority} & \hd{Why} \\
\midrule
Periodic re-zero during a streaming run & high &
  The 25\% PL droop (\S\ref{sec:droop}) walks the operating point off the
  calibrated flank over 30\,s. Currently only detrended in post. \\
Root-cause the burst replay & high &
  Would recover $\sim$15\% of burst rate and remove the need for the drop policy. \\
Real-time mains notch in the acquisition path & medium &
  The current notch is offline, non-causal, and assumes stationarity. At $\ge$1\,kHz a
  causal notch is feasible and the line is 10$\times$ the floor. \\
Fold the multipoint notebook onto the same modules & medium &
  \file{twopoint_spectra.py} and \file{twopoint_postprocess.py} are mode-generic;
  the multipoint path still has its own analysis code. \\
Calibrated field reference & medium &
  Every $\eta$ here is a readout-chain figure (\S\ref{sec:notaclaim}). A known
  applied field would turn them into system claims. \\
\co{remove_offset} is documented but not implemented & low &
  \co{NVAveragerProgram.acquire()} accepts the argument and never uses it. Harmless
  today (both paths are consistent) but a latent trap. \\
Allan deviation on the long runs & low &
  Needed for the drift/averaging-time story; the June time-integration data exists
  but has not been analysed drift-aware. \\
\bottomrule
\end{tabular}
\end{center}

\section{Recommended operating point today}

\begin{keybox}[Until S1 resolves the streaming question]
\begin{description}[leftmargin=1.4em,style=nextline]
  \item[For best sensitivity --- Method B, burst.] $\tau=\SI{120}{\micro\second}$,
  \co{BURST_REPS = 500}, \co{BURST_ON_STALE = "drop"}. 38\,\etaunit\ at an effective
  3.1\,kHz. Accept the $\sim$15\,ms gap every 135\,ms.
  \item[For a continuous trace --- Method C, streaming.] $\tau=\SI{120}{\micro\second}$,
  \co{STREAM_TARGET_HZ = 1000}. 1041.7\,Hz, gapless, exact timestamps, but
  $\sim$6$\times$ the noise until \S\ref{sec:streamnoise} is closed.
  \item[For slow work and sanity checks --- Method A, averaged.]
  $\tau=\SI{120}{\micro\second}$, \co{AVG_REPS_PER_BATCH = None} so the call floor is
  filled. $\sim$88\,Hz, no open defects --- but remember 60\,Hz is aliased into the
  band here.
\end{description}
\end{keybox}


\chapter{Corrections register}
\label{ch:retractions}

Every claim from an earlier document that this one contradicts, in one place.

\begin{center}
\small
\begin{longtable}{@{}p{3.2cm}p{4.6cm}p{4.6cm}p{2.4cm}@{}}
\toprule
\hd{Source} & \hd{Earlier claim} & \hd{Current position} & \hd{Evidence} \\
\midrule
\endhead
2026-08-06 doc \S1, \S2 &
  ``The rate is FPGA-bound, not host-bound''; host round-trip $\sim$1.6\,ms &
  \textbf{Retracted.} Host-bound, $\sim$10--11.4\,ms per call; the FPGA runs
  underneath it &
  \S\ref{sec:floorevidence} \\[3pt]
2026-08-06 doc \S2.6 &
  $\tau$ is ``the main lever''; \co{reps} trades against rate one for one &
  \textbf{Both backwards.} Neither affects the rate below the floor; \co{reps} is a
  free \emph{sensitivity} lever &
  \S\ref{sec:knobs} \\[3pt]
2026-08-06 doc \S2.4 &
  ``Host round-trip 1.60\,ms (14.1\%)'' as a budget row &
  That was a \emph{residual} under an assumed serial model, not a measurement &
  \S\ref{sec:floorevidence} \\[3pt]
2026-08-06 doc \S4.3 &
  \co{config_all} / never-stopped tProc is the root cause of the burst replay &
  \textbf{Disproved.} The fix was applied and 49.9\% still replayed &
  \S\ref{sec:replay} \\[3pt]
2026-08-06 doc \S3.2 &
  ``No 60\,Hz line, no discrete tones'' &
  The line is there and is 10$\times$ the floor; at 87.5\,Hz it is \emph{aliased} to
  $\sim$28\,Hz and the burst runs that could have resolved it were half replays &
  \S\ref{sec:alias} \\[3pt]
2026-08-06 doc, 1st draft &
  $\eta \approx 19$\,\etaunit\ per rep &
  185\,\etaunit. The draft paired a reps-averaged $\sigma$ with a single-rep time,
  omitting $\sqrt{\text{reps}}\approx 10$ &
  \S\ref{sec:conventions} \\[3pt]
2026-08-06 doc, $\eta$ convention &
  $\eta = \sigma_B\sqrt{\Delta t}$ &
  $\eta = \sigma_B\sqrt{2\Delta t}$, matching what the notebook prints; a
  $\sqrt{2}$ difference &
  \S\ref{sec:conventions} \\[3pt]
2026-08-14 doc, 1st draft &
  Streaming excess is ``4.2$\times$'' &
  5.3--6.3$\times$. The 4.2 used the \emph{raw} burst floor; the cleaned floor is
  38.6 rather than 55\,\etaunit &
  \S\ref{sec:streamnoise} \\[3pt]
Notebook Step 4 comment, pre-2026-08-12 &
  ``$\sim$10\,ms of fixed host overhead, FPGA work that fits inside it is FREE'' &
  \textbf{Substantially correct}, and wrongly ``corrected'' on 2026-08-12. Restored &
  \S\ref{sec:model} \\[3pt]
\co{time_per_rep()}, pre-2026-08-12 &
  $n_\text{slots}\times(\tau + 2\,t_\text{relax})$ &
  $n_\text{slots}\times\tau$; the relax windows are absorbed &
  \S\ref{sec:asm} \\[3pt]
Burst cadence, pre-2026-08-16 &
  From \co{acq_seconds/n_samples}: \SI{269}{\micro\second} &
  \SI{240.0}{\micro\second} exactly, from the readout quantum. The old value
  included the host share &
  \S\ref{sec:quantum} \\[3pt]
Despiking, pre-2026-08-16 &
  Filter the two parked channels independently &
  Filter the \emph{differential}. Per-channel raises $\sigma$ &
  \S\ref{sec:despikebug} \\
\bottomrule
\end{longtable}
\end{center}

\begin{keybox}[The methodological lesson]
Four of the twelve rows above come from the same mistake: \textbf{fitting a model
with $n$ free parameters to data taken at one operating point}. The 2026-08-06
session ran at a single $(\tau, \text{reps})$, and every timing conclusion drawn
from it was under-determined --- the serial-versus-overlapped question simply cannot
be answered from one point, however carefully the residuals are computed.

The 2026-08-14 session settled it in an afternoon purely by varying two knobs. Any
future timing claim should span at least three operating points before it is
written down.
\end{keybox}


\chapter{Reproducing everything}

\section{Offline, no hardware}

\begin{quote}\small\ttfamily
\# every number and figure in this document\\
python scripts/analyze\_twopoint\_0814.py\\[4pt]
\# post-process any run; the mode is detected from the columns\\
python -m twopoint\_postprocess <csv> [-{}-outdir DIR] [-{}-notch-mains]\\[4pt]
\# burst-specific CLI with the original flags\\
python scripts/clean\_burst\_lockin.py <live\_csv> -{}-outdir analysis/\\[4pt]
\# the 2026-08-06 analysis, unchanged\\
python scripts/analyze\_twopoint\_session.py\\[4pt]
\# static check of the config\_all signature\\
python scripts/profile\_twopoint\_acquire.py -{}-check-only\\[4pt]
\# rebuild this document\\
cd docs/twopoint\_master\_reference \&\& latexmk -pdf TWOPOINT\_MASTER\_REFERENCE.tex
\end{quote}

\section{On the rig}

\begin{quote}\small\ttfamily
python scripts/profile\_twopoint\_acquire.py -{}-floor -{}-tau 120\\
python scripts/profile\_twopoint\_acquire.py -{}-compare-read-paths -{}-tau 120 -{}-reps 500\\
python scripts/profile\_twopoint\_acquire.py -{}-stream-scan\\
python scripts/profile\_twopoint\_acquire.py -{}-ip 192.168.0.103\\
\hspace*{2em}\# the full per-RPC profile (fills section~\ref{sec:rpc})
\end{quote}

\section{Data used by this document}

\begin{center}
\small
\begin{tabular}{@{}p{6.4cm}p{8.0cm}@{}}
\toprule
\hd{Path} & \hd{Contents} \\
\midrule
\file{data/results/} \file{081426 (Sensitivity increase update)/} &
  2026-08-14: 4 averaged + 1 stream + 1 burst two-point run, all at
  $\tau$=\SI{120}{\micro\second}; 12-point integration sweep; 8-peak reference sweep;
  multipoint run. 101\,MB. \\
\file{data/twopoint_lockin/08062026/} &
  2026-08-06: drone sweeps, 4 burst runs, 12 averaged runs, 17-point integration
  sweep (92 ODMR sweeps). \\
\file{docs/2026-08-14_twopoint_methods/} \file{tables/} &
  Generated: \co{timing.csv}, \co{mode_budgets.csv}, \co{modes.csv},
  \co{read_path.csv}, \co{per_rep_noise.csv}, \co{spectra.csv}. \\
\bottomrule
\end{tabular}
\end{center}
